\section{An elementary bound on representation degrees}
\label{sec:degrees}

Proposition~\ref{prop:lift:orbit-sum} reduced the Fourier estimate to a
finite sum of squared degrees of small representations of each block
group. We now bound that sum. For integers $s,D,k\ge1$ and a
group $H\cong S_s$, we have
\begin{align}
 \sum_{\substack{\pi\in\widehat H\\\dim\pi\le D^{1/k}}}(\dim\pi)^2
 &=\sum_{\substack{\pi\in\widehat H\\\dim\pi\le D^{1/k}}}
       (\dim\pi)^{18}(\dim\pi)^{-16}\notag\\
 &\le D^{18/k}\sum_{\pi\in\widehat H}(\dim\pi)^{-16}.
\label{eq:degrees:truncated-sum}
\end{align}
Thus a uniform bound on the sum of inverse sixteenth powers of degrees will
control every block size in the preceding proposition. We also prove that
the sum tends to zero after the trivial and sign representations are
removed. That second statement will control the full-group Fourier sum
in the completion.

Recall that the complex irreducible representations of $S_s$ are indexed
by partitions $\lambda\vdash s$, and that their degrees $f^\lambda$ count
standard Young tableaux of shape $\lambda$. Such a tableau fills the
diagram with $1,\ldots,s$, increasing along every row and column. The
partitions $(s)$ and $(1^s)$ afford the trivial and sign representations,
respectively; transposing a tableau gives
$f^{\lambda'}=f^\lambda$.
These are the classical Specht-module and standard-tableau facts
\cite{sagan,etingof}. Liebeck and Shalev prove the stronger estimate
$\sum_{\lambda\vdash s}(f^\lambda)^{-u}=2+O_u(s^{-u})$
as $s\to\infty$, for every fixed $u>0$
\cite[Theorems~1.1 and~2.6]{liebeck-shalev2004}.
The following elementary proof at exponent $16$ keeps this application
independent of that stronger estimate; it uses only the tableau facts just
recalled.

\begin{lemma}[Summability of inverse degrees]
\label{lem:degrees}
There is an absolute finite constant
\begin{equation}
  C_0:=\sup_{s\ge 1}\sum_{\lambda\vdash s}(f^\lambda)^{-16}<\infty.
  \label{eq:degrees:uniform}
\end{equation}
Moreover,
\begin{equation}
  \sum_{\substack{\lambda\vdash s\\
                  \lambda\notin\{(s),(1^s)\}}}
       (f^\lambda)^{-16}\longrightarrow 0
       \qquad\text{as }s\longrightarrow\infty.
  \label{eq:degrees:vanishing}
\end{equation}
The excluded partitions are interpreted as a set, so when $s=1$ only
the single partition $(1)$ is excluded.
\end{lemma}

\begin{proof}
We first record two elementary counting observations.
If a diagram has first row of length $a$ and exactly $b$ boxes below it,
where $0\le b\le a$, then
\begin{equation}
  f^\lambda\ge \binom{a}{b}.
  \label{eq:degrees:tableau-construction}
\end{equation}
This is the tableau construction of Liebeck and Shalev
\cite[Lemma~2.1]{liebeck-shalev2004}, with their parameters $n=a+b$ and
$k=b$. To see it directly, place $1,\ldots,b$ in the first $b$ boxes of
the first row.
Choose any $b$ labels from $\{b+1,\ldots,a+b\}$ for the boxes below
that row, and put them in increasing row-major order: read rows from
top to bottom, and each row from left to right. Fill the rest of the
first row increasingly with the remaining labels. Every box below
the first row lies in one of its first $b$ columns. Its top-row
predecessor therefore has a label at most $b$, while all lower labels
are larger than $b$. Row-major filling respects both row and column
inequalities within the lower diagram, so this is a standard tableau.
Different choices of lower labels give different tableaux. When $b=0$
the construction is the unique tableau of a single row.

Also, if $\mu$ is a Young subdiagram of $\lambda$, every standard
tableau of $\mu$ extends to one of $\lambda$: retain its labels and
fill the remaining boxes in row-major order with the successive
larger labels. Because a Young subdiagram is closed under moving up
and left, all required predecessors have already been filled.
Consequently,
\begin{equation}
  f^\lambda\ge f^\mu.
  \label{eq:degrees:extension}
\end{equation}

Write $p(t)$ for the number of partitions of $t$. For $t\ge 1$,
regarding a partition as an ordered composition gives
\begin{equation}
  p(t)\le 2^{t-1}\le 2^t,
  \label{eq:degrees:partition-count}
\end{equation}
since a composition is specified by cuts in the $t-1$ gaps between
$t$ successive units. We set $p(0)=1$.

For a partition $\lambda\vdash s$, put
\[
  L=\max\{\lambda_1,\lambda'_1\},\qquad h=s-L.
\]
After transposing if necessary, assume its first row has length $L$.
This does not change its degree. The case $h=0$ consists precisely of
the row and column partitions, which are excluded from
\eqref{eq:degrees:vanishing}. We divide the other diagrams into three
classes.

\smallskip
\noindent\textit{A long first row or column: $1\le h\le s/3$.}
For fixed $h$, there are at most $2p(h)$ original shapes: after a
possible transpose, the diagram below the first row is a partition
of $h$. By \eqref{eq:degrees:tableau-construction},
\[
  f^\lambda\ge\binom{s-h}{h}
  \ge \left(\frac{s-h}{h}\right)^h\ge 2^h.
\]
Here the middle inequality follows by writing
$\binom{a}{b}=\prod_{j=0}^{b-1}(a-j)/(b-j)$ and bounding each factor
below by $a/b$ when $a\ge b\ge1$.
Thus the total inverse-degree contribution with a given $h$ is at most
\[
  2\cdot 2^h\binom{s-h}{h}^{-16}
  \le 2\cdot 2^{-15h}.
\]
For each fixed positive $h$, the binomial coefficient tends to infinity
with $s$. The displayed geometric majorant is summable over $h$.
Extending a contribution by zero when $h>s/3$, dominated convergence
for this series shows that the total contribution of this class tends
to zero.

\smallskip
\noindent\textit{Intermediate width: $h>s/3$ and $L\ge s/8$.}
Here $L<2s/3$, so $h>L/2$. Set $b=\lfloor L/2\rfloor$ and take
the subdiagram $\mu$ consisting of the whole first row and the first
$b$ boxes below it in row-major order. There are enough boxes because
$h\ge b$, and an initial segment of row-major order in the lower
diagram is again a Young diagram. Applying
\eqref{eq:degrees:extension} and
\eqref{eq:degrees:tableau-construction} gives
\[
  f^\lambda\ge \binom{L}{\lfloor L/2\rfloor}
  \ge \frac{2^L}{L+1}
  \ge \frac{2^{s/8}}{s+1}.
\]
The central binomial coefficient is the largest of the $L+1$
coefficients whose sum is $2^L$. Using
\eqref{eq:degrees:partition-count}, the total contribution of this
class is therefore at most
\begin{equation}
  2^s\left(\frac{s+1}{2^{s/8}}\right)^{16}
  =(s+1)^{16}2^{-s}\longrightarrow0.
  \label{eq:degrees:intermediate}
\end{equation}

\smallskip
\noindent\textit{Small width and height: $L<s/8$.}
Number row and column indices starting at $1$, and group boxes by
their index sum. There are at most $2L-1$ nonempty such diagonals.
Fill the diagonals in increasing index-sum order, assigning the next
available consecutive labels to each diagonal in any order. Every
row or column predecessor has a strictly smaller index sum, so every
such filling is a standard tableau. If the nonempty diagonal sizes
are $u_1,\ldots,u_q$, then
\[
  f^\lambda\ge \prod_{j=1}^q u_j!
  \ge \prod_{j=1}^q 2^{u_j-1}
  =2^{s-q}\ge 2^{3s/4}.
\]
We used $u!\ge2^{u-1}$ for integers $u\ge1$, and
$q\le2L-1<s/4$. The inverse-degree contribution of this class is
at most
\begin{equation}
  2^s\bigl(2^{3s/4}\bigr)^{-16}
  =2^{-11s}\longrightarrow0.
  \label{eq:degrees:small-width}
\end{equation}

These three classes cover all partitions except the row and column,
so they prove \eqref{eq:degrees:vanishing}. For every $s\ge2$ those
two partitions are distinct and each has degree $1$. The complete
sum in \eqref{eq:degrees:uniform} consequently tends to $2$ as
$s\to\infty$. For each fixed positive $s$ it is a finite sum of
positive terms, since there are finitely many partitions and every
degree is a positive integer. The tail of this convergent sequence
is bounded, and its finitely many initial values have a finite
maximum. This proves \eqref{eq:degrees:uniform}, including $s=1$,
where the complete sum is $1$.
\end{proof}

In representation notation, the uniform conclusion of the lemma is
\begin{equation}
\label{eq:lift:degree-constant}
 C_0=\sup_{s\ge1}\sum_{\pi\in\widehat{S_s}}(\dim\pi)^{-16}<\infty.
\end{equation}
We can now finish the transfer from bounded coset densities to Fourier
decay. The bound is uniform over the block sizes because $C_0$ is uniform
over all positive integers $s$.

\begin{corollary}[From bounded coset densities to Fourier decay]
\label{prop:lift}
Let a finite set $X$ be partitioned into $k>18$ nonempty sets
$M_1,\ldots,M_k$, and put $G=\Sym(X)$. Let $H_a$ be the subgroup that
permutes $M_a$ and fixes $X\setminus M_a$ pointwise. Suppose that a
probability measure $\nu$ on $G$ satisfies, for some $B\ge1$,
\[
 \nu(gH_a)\le B\frac{|H_a|}{|G|}
 \qquad(g\in G,\ 1\le a\le k).
\]
Then every irreducible complex unitary representation $\rho$ of $G$, of
dimension $D$, satisfies
\begin{equation}
\label{eq:lift:general-fourier}
 \|\widehat\nu(\rho)\|_{\op}
 \le\sqrt{kBC_0}\,D^{-(1-18/k)/2}.
\end{equation}
Here $C_0$ is the absolute constant of Lemma~\ref{lem:degrees}. The block
sizes need not be equal.
\end{corollary}

\begin{proof}
For every $a$, the group $H_a$ is isomorphic to $S_{|M_a|}$. Hence
\eqref{eq:degrees:truncated-sum} and Lemma~\ref{lem:degrees} give
\[
 \sum_{\substack{\pi\in\widehat H_a\\\dim\pi\le D^{1/k}}}(\dim\pi)^2
 \le C_0D^{18/k}.
\]
In the orbit construction of Proposition~\ref{prop:lift:orbit-sum}, this
recovers $\dim W_a\le C_0D^{18/k}$, and the component bound
\eqref{eq:lift:component-bound} becomes
\[
 |\langle z,\widehat\nu(\rho)u_a\rangle|^2
 \le BC_0D^{-1+18/k}\|u_a\|^2\|z\|^2.
\]
Substituting the displayed degree bound into
\eqref{eq:lift:orbit-sum} yields
\[
 \|\widehat\nu(\rho)\|_{\op}^2
 \le kBC_0D^{-1+18/k}.
\]
Taking square roots proves \eqref{eq:lift:general-fourier}.
\end{proof}
